% Transcription — fonds Alexandre Grothendieck, shelfmark 138, batch 3 (pages 41-60).
% Established from the facsimile archives/batches/138/batch-03.pdf (cover sheet
% excluded; batch page 1 = archivists' page 41), per the skill
% transcribe-grothendieck.
%
% Pass: Opus 5.5 (claude-opus-5-5), 2026-09-30 — first pass, unchecked against the pages by a human.
%
% Thirteen of the twenty pages are transcribed. Pages 42, 44, 47, 49, 51 and 55
% are unrelated administrative or typed versos with nothing in his hand; they
% are absent. Page 56 is a cover sheet bearing only his words « autom. des
% groupoïdes », used as the heading of pages 57-60 and given no \page.
%
% Contents, in outline:
%   Page 41 — the groups pi(I) = Zhat^I / Zhat, pi_n(I), an extension E''_1(I)
%     by S_I x Zhat^*, the sets E(I,i), E(I,omega) for I a 3-element set.
%   Page 43 — genus of the Fermat curve C_n; C_2 : x^2+y^2+z^2 = 0 read as a
%     weighted triangulation: the octahedron, automorphism group S_4 as a
%     semidirect product of S_3 by (Z_2)^2; its three quotients of order 2.
%   Pages 45-46 — X = C_2, its ramification divisor, the pro-covering X_inf;
%     the local picture over 0, 1, inf of x_0 = 1/2 - 1/4 (z + z^-1);
%     D_inf = (Z/2) . mu_inf; sketches of ramification data.
%   Page 48 — the map (x,y,z) -> (x(...), y(...), z^3) on the conic;
%     a Lagrange-type computation p - (p'/q') q = 0.
%   Page 50 — rational parametrisation of the conic, the squaring map
%     (x^2 - y^2, 2xy, x^2 + y^2), maps X -> X' of degrees 2, n.
%   Pages 52-53 — a diagram of coverings X_1 <- X_2 <- X_4 ... with Galois
%     groups mu_n^3 / mu_n.
%   Page 54 — X + iY = ((y - ix)/z)^2 and its powers.
%   Pages 57-60 — « Automorphismes de groupoïdes »: functors from a groupoid
%     over a given map on objects, criteria for faithful, fully faithful,
%     equivalence, isomorphism; the exact sequence 1 -> Aut^0(C) -> Aut(C) ->
%     S_{C_0} -> 1. Page 60 breaks off mid-sentence.
%
\documentclass[11pt,a4paper]{article}
\input{../preamble/grothendieck}

\folder{138}
\batch{3}
\pages{41}{60}
\foldertitle{Cartes. Etude arithmétique (1976, 77 ?) : bon de commande (s.d.), notes manuscrites (s.d.).}
\dating{[vers 1976-1978]}
\watermark{Édition de démonstration}

\begin{document}

\page{41}
\note{la page s'ouvre sur des formules sans phrase d'introduction ; ce qui les précède se trouve avant ce lot.}

$I \in \mathrm{Ob}\,\mathrm{Ens}_3$

\[
  \pi(I) = \hat{\mathbb{Z}}^{I} / \hat{\mathbb{Z}},
  \qquad
  \pi_n(I) = (\mathbb{Z}/n)^{I} / (\mathbb{Z}/n)
\]

\marginal{à droite, séparé par un trait vertical : $\mathfrak{S}_I \times \hat{\mathbb{Z}}^{\times}$, avec deux flèches vers « \uncertain{structurel} » (sur $\mathfrak{S}_I$) et « \uncertain{hom.} » (sur $\hat{\mathbb{Z}}^{\times}$), et « y opère »}

\[
  \mathcal{E}''_{1}(I) \simeq \bigl(\mathfrak{S}_I \times \hat{\mathbb{Z}}^{\times}\bigr) \cdot_{1/2} \pi(I)
\]
\note{l'indice $\tfrac12$ sous le point est celui qu'il met aussi p.~46 ($\mathbb{D}_\infty = (\mathbb{Z}/2)\cdot_{1/2}\mu_\infty$) ; il semble noter un produit semi-direct (« $\tfrac12$ direct », p.~43).}

\[
  E(I, i) \simeq \hat{\mathbb{Z}}^{(\pm 1)}_{\wedge}\bigl(I \setminus \{i\}\bigr)
  \qquad (i \in I)
\]
\[
  E(I; I) \simeq \coprod_{i \in I} E(I, i)
\]

\struck{\ill{}} \quad
$\omega \simeq \mathrm{Or}(I) \wedge \{1, \bar{\jmath}\} \simeq \mathrm{Or}(I) \subset \pi_3 = (\mathbb{Z}/3)^{I}/(\mathbb{Z}/3) = \mathbb{F}_3^{I}/\mathbb{F}_3$
\note{le second élément de l'accolade, lu $\bar{\jmath}$, est douteux.}

\[
  E(I, \omega) = p_3^{-1}(\omega) \subset \pi_\infty(I)
\]
\[
  0 \to \pi_\infty(I) \to \pi_\infty(I) \to \pi_3(I) \to 0,
  \qquad \omega \in \pi_3(I)
\]
\note{le premier terme de la suite, lu $\pi_\infty(I)$, est surchargé et peut-être biffé : lecture \uncertain{douteuse}.}

Une flèche descend de \struck{\ill{}} vers $E(I, i)$.

\page{43}
Genre de la courbe de Fermat $C_n$ (qui est une courbe \struck{\ill{}} \add{\ill{}} lisse de degré $n$)
\[
  g_n = \frac{(n-1)(n-2)}{2}.
\]
Le genre est $0$ ssi $n = 1, 2$, il est $1$ ssi $n = 3$. Pour $n = 4$ il est $3$, pour $n = 6$ il est $10$, pour $n = 12$ il est $55$ …

Considérons le cas $n = 2$
\[
  C_2 : x^2 + y^2 + z^2 = 0
\]
comme une triangulation \uncertain{pondérée}. Elle est formée des $4 \cdot 2 = 8$ triangles, chaque type \ill{} \uncertain{intervenant} $2$ fois, avec degré \ill{} $4$ \uncertain{(p.~2.2)} \ill{} pour les \ill{} \ill{}. Sauf erreur, c'est l'\emph{octaèdre} \add{sphérique}, dont le groupe des automorphismes orientés est bien $\simeq \mathfrak{S}_4$ d'ordre $24$, et peut s'écrire comme produit $\tfrac12$ direct de $\mathfrak{S}_3$ par $(\mathbb{Z}_2)^2$ !

\emph{Divisant} par les trois s-groupes $\mathbb{F}_2$ de $\pi_2 \simeq \mu_2^3/\mu_2$, on trouve les $3$ revêtements d'ordre $2$ de $C_1 \simeq \mathbb{P}^1_{\mathbb{Q}}$ qui correspondent au \struck{dièdre} \add{\uncertain{monoèdre}} \struck{\ill{}} (\uncertain{cardinal}) et \uncertain{à ses} transformés par $\mathfrak{S}_3$.

\drawing{à gauche, une sphère parcourue de trois grands cercles, les points d'intersection marqués $0$, $1$, $\infty$ (chacun deux fois) ; au-dessous, une seconde figure : une courbe fermée contenant un ovale, avec les points $\infty$, $0$, $1$, $0$, $\infty$ et deux traits marqués $1$.}

\note{la note suivante est encadrée, en bas à gauche de la page, à côté de la seconde figure ; sa prose se lit mal.}
\emph{NB} L'octaèdre ou subdivision \uncertain{barycentrique} du \struck{\ill{}} \uncertain{biédre} (qui est \uncertain{autodual}) \struck{\ill{}} et a le type de ramification
\[
  \begin{matrix} 2 & 2 & 2 \\ 2 & 2 & 2 \end{matrix}
\]
\uncertain{en particulier} de celle du \struck{\ill{}} \uncertain{dièdre}
\[
  \left(\begin{matrix} 2 & 2 & n \\ 2 & 2 & n \end{matrix} \;\text{pour } n = 2\right),
\]
ce qui \emph{confirme}, via la \uncertain{unicité} de la \ill{} \uncertain{rev.-théorie} \ill{} \uncertain{à conjugaison} \uncertain{près} par $2$ …

\page{45}
Posons pour abréger $X = C_2$, soit $D$ le diviseur de ramification sur $X$ (de degré $6$) --- \uncertain{son corps} de \uncertain{définition} est $\mathbb{Q}(i)$ ; \add{posons $X^{*} = X \setminus D$} \struck{posons $X^{*} = X \setminus D$},
\struck{et soit $\widetilde{X^{*}}$ le rev. universel abélien de $X^{*}$}

Soit $X_\infty$ le \add{pro-}revêtement \struck{abélien} \add{principal} de $X$ défini par la famille des $X_n$ ($n$ pair), il a comme groupe $(\pi_\infty)_{\struck{X}\,2}$
\note{l'indice de $(\pi_\infty)$ est surchargé : un $X$ barré, puis un chiffre, peut-être $2$. La phrase s'arrête là ; le reste de la page est blanc.}

\page{46}
\[
  \begin{array}{ll}
    X_n = \widehat{\mathrm{Spec}}\, k[x_n] & x_n^{\,n} = x_1 \\
    \downarrow & \\
    X_1 = \widehat{\mathrm{Spec}}\, k[x_1 = z] & \\
    \downarrow & \\
    X_0 = \widehat{\mathrm{Spec}}\, k[x_0] &
      x_0 = \frac12 - \frac14\bigl(z + z^{-1}\bigr)
  \end{array}
\]
\drawing{en marge gauche de la première flèche, une étoile à branches (le revêtement $x_n^n = x_1$) ; en marge de la seconde, un cercle marqué $-1$ et $1$.}

Pour $X_n \to X_1$ (encadré : $n \{ \vdots\; n$) : ramifications en dessus des pts $0, \infty$ (où la ramification est $n$) ;
\note{ligne ajoutée au-dessous :} \add{pts de ramification en haut : $0, \infty$}

Pour $X_1 \to X_0$ (encadré : $2\;2 \mid 1\;1$) : ramification au dessus de
\begin{itemize}
  \item $0$ \add{($z = 1$, ramifié $2$)} \struck{$\infty$ \ill{} (pts \ill{} dessus)} ;
  \item $1$ ($z = -1$, ramifié $2$) \struck{\uncertain{idem}} ;
  \item $\infty$ ($0$ et $\infty$, ramification $1$).
\end{itemize}

\[
  \mathbb{D}_\infty = (\mathbb{Z}/2) \cdot_{1/2} \mu_\infty
\]

\drawing{au milieu de la page, un schéma des fibres : au-dessus de $0$ et de $\infty$ en bas, deux flèches venant de « $0\,(1)$ » et « $\infty\,(1)$ », elles-mêmes images de « $0\,(n)$ » et « $\infty\,(n)$ » ; à côté, « $1\,(1)$ » (entouré : « $n$ fois ») $\to$ « $1\,(2)$ » $\to 0$, et, en partie biffé, « $-1\,(2)$ » $\to 1$ ; à droite, « $n$ pair ». À gauche, deux lignes de points : $0\,(2)$, $1\,(2)$, $\infty_1\,(1)$, $\infty_2\,(1)$, puis $0$, $1$, $\infty$.}

\drawing{en bas, un ovale portant alternativement des traits marqués $0$ et des points marqués $1$ (quatre de chaque), avec $\infty$ au-dessus et au-dessous de l'ovale.}

\page{48}
\[
  z' = z^3, \qquad
  y' = y\,(y^2 - 3x^2), \qquad
  x' = -x\,(x^2 - 3y^2)
\]
\begin{align*}
  & z^6 + y^2\,(y^4 + 9x^4 - 6x^2y^2) + x^2\,(x^4 + 9y^4 - 6x^2y^2) \\
  &\quad = x^6 + y^6 + z^6 + 3\,x^2y^2\,(x^2 + y^2) = \text{\struck{$(x^2 + y^2 + z^2)^3$}} \\
  &\quad = x^6 + y^6 + z^6 - 3\,x^2y^2z^2
\end{align*}
\note{à côté de la dernière ligne, un début non achevé : $x^6 + y^6 + z^6 + \cdots$. Les signes de la seconde ligne sont écrits vite ; le premier $+$ entre $x^6$ et $y^6$ pourrait être un $-$.}

\[
  \Bigl(\frac{y}{z}\Bigr)^2 + \Bigl(\frac{x}{z}\Bigr)^2 = -1
  \qquad a + ib \qquad \text{\struck{$a^2 + b^2 = -1$}}
\]
\struck{$x^2 + y^2 + z^2$}

\drawing{deux copies de la droite, $X$ et $X'$, avec les points $0$, $1$, $-1$, $2$, $\infty$ sur $X$ et $0$, $1$, $\infty$ sur $X'$ ; des flèches courbes envoient $0 \mapsto 0$, $1$ et $-1 \mapsto 1$, $\infty \mapsto \infty$ ; deux traits verticaux, le second marqué $n$.}

\begin{gather*}
  \frac{p(x)}{q(x)} = \lambda, \qquad p(x) - \lambda\, q(x) = 0, \qquad p'(x) - \lambda\, q'(x) = 0 \\
  \lambda = \frac{p'}{q'}, \qquad p - \frac{p'}{q'}\, q = 0, \qquad \text{\struck{$q'p -$}}\; p'q - q'p = 0
\end{gather*}
$(x^2)^2 = (\pm x \cdots)^2$ \ill{} \qquad
$\dfrac{az^2 + bz + c}{z^2}$ \qquad
\ill{} $2z\,(az^2 + bz + c) - \cdots$
\note{dans cette dernière ligne, le $\pm x$ et le coefficient $a$ (surchargé) sont des lectures \uncertain{douteuses}.}

\page{50}
\[
  \cos 2u = \frac{1 - \mathrm{tg}^2 u}{1 + \mathrm{tg}^2 u},
  \qquad
  \sin 2u = \frac{2\,\mathrm{tg}\, u}{1 + \mathrm{tg}^2 u}
\]
\[
  2 \sin u \cos u = 2 \cos^2 u \;\mathrm{tg}\, u =
\]
\[
  \left\{
  \begin{aligned}
    x &= \frac{1 - t^2}{1 + t^2} \\
    y &= \frac{2t}{1 + t^2}
  \end{aligned}
  \right.
  \qquad t = y/x \quad \text{quand } z = 1
\]
\note{« quand » est une lecture \uncertain{douteuse}.}
\note{sous $t = y/x$, une expression biffée et surchargée, commençant par $1 + t + \cdots$.}
\[
  x = 1 - t^2, \qquad y = 2t, \qquad z = 1 + t^2
\]
\struck{$y \quad t = \pm i$} \qquad $t = 0$

\[
  z = 0 \text{ i.e. } t = \mp i \;\Longrightarrow\; x = 2,\ y = \mp 2i \;\sim\; x = 1,\ y = \mp i
\]
\[
  \sim\; x = \pm i,\ y = 1
\]
\[
  y = 0 \text{ i.e. } t = 0 \qquad (1, 0, 1)
\]

\[
  x^2 + y^2 + z^2 = 0 \qquad C_2 = X \longrightarrow \mathbb{P}^1 = \widehat{\mathrm{Spec}}\,(k[t])
\]
\note{le « $\widehat{\mathrm{Spec}}$ » est \uncertain{douteux} : le mot sous le chapeau est à peine formé.}
\[
  (\pm i, 1, 0) \longrightarrow t = \mp i, \qquad
  (\pm i, 0, 1) \longmapsto t = 0, \qquad
  t = \frac{y}{x}
\]
\[
  x' = 1 - t^2 = 1 - \frac{y^2}{x^2} = \frac{x^2 - y^2}{x^2}, \qquad
  y' = 2t = \frac{2y}{x}, \qquad
  z' = 1 + t^2 = \frac{x^2 + y^2}{x^2}
\]
\[
  \sim\; (x', y', z') = (x^2 - y^2,\ 2xy,\ x^2 + y^2)
\]

\[
  \left\{
  \begin{aligned}
    x' &= x^2 - y^2 \\
    y' &= 2xy \\
    z' &= x^2 + y^2
  \end{aligned}
  \right.
  \qquad
  \begin{aligned}
    x'^2 + y'^2 - z'^2 &= x^4 + y^4 - 2x^2y^2 + 4x^2y^2 \\
    &\quad - (x^2 + y^2)^2 = 0
  \end{aligned}
\]

\begin{tikzcd}
  X \arrow[r, "2"] \arrow[dr, bend right=20, "f_n"'] & X' \arrow[r, "?"] & X' \\
  & X \arrow[ur, "2"'] &
\end{tikzcd}
\note{le label de la flèche $X' \to X'$, noté « ? », est \ill{} : il se lit « $n$-\ill{} », peut-être « $n$-ifian » ; la flèche courbe est marquée $f_n$.}

\[
  s = -\frac{t + i}{t - i} = -\frac{(t + i)^2}{t^2 + 1}
  = -\frac{t^2 - 1}{t^2 + 1} - \frac{2it}{t^2 + 1}
\]
\[
  \Bigl(\frac{t + i}{t - 1}\Bigr)^{n} =
\]
\note{au dénominateur de la dernière formule il écrit $t - 1$, là où la ligne précédente a $t - i$. Un signe \ill{} précède $s$ dans la marge.}

\page{52}
\[
  \pi_n = \mu_n^{3} / \mu_n
\]
\[
  \widetilde{X}_2(p, q, r) = X_2(2p, q, 2r)
\]
\marginal{en haut à droite : $0' \mapsto{}$ \uncertain{$\infty''$}, $1' \mapsto 1''$, $\infty' \mapsto{}$ \uncertain{$0''$} ; au-dessous, soulignés deux à deux : $0' \mapsto 1'$, $1' \mapsto \infty'$ ?}

\drawing{un diagramme de revêtements, trop surchargé pour être mis en \texttt{tikz-cd} sans choisir entre des flèches biffées. On y lit :
$X_2(p,q,r) \to X_2$ (étiquette $\mu_p \times \mu_q \times \mu_r$) ;
$\widetilde{X}_2(p,q,r) \to \widetilde{X}_2$ (étiquette biffée, puis $\mu_{2p} \times \mu_{2q} \times \mu_{2r} \times \pi_n$ ?) ;
$X_2(2,2,2) = \widetilde{X}_2 \to X_2$ (étiquette $\mu_2^{3}$) et $\widetilde{X}_2 \to X_4$ (étiquette $\mu_2$) ;
$\widetilde{X}_2(n) \to \widetilde{X}_2$ (étiquette $\pi_n$) et $\widetilde{X}_2(n) \to X_{4n}$ (étiquette $\mu_{2n}$ ?) ;
$X_{4n} \to X_4$ (étiquette $\pi_4$) ;
$X_4 \to X_2$ (étiquette $\pi_2$), $X_4 \to X_1$ (étiquette $\pi_4$) ;
$X_2 \to X_1$ (étiquette $\pi_2$) ;
une longue flèche de la droite vers $X_2$ sous une étiquette très biffée ($\mu_2^{3}$ ?), et une flèche courbe de $X_1$ vers la droite, étiquetée $\pi_{2n}$ ;
un grand arc part du bas et remonte vers $\widetilde{X}_2(p,q,r)$.}

\[
  \mu_{2p} \times \mu_{2q} \times \mu_{2r} \times \pi_n
\]
si $p = q = r = 1$, $n = 2$, on trouve
\[
  \mu_2 \times \mu_2 \times \mu_2 \times \pi_2 \twoheadrightarrow \mu_2^{3}/\mu_2
\]
\note{la flèche finale est un trait souligné sous $\pi_2$ ; le sens « $\twoheadrightarrow$ » est une lecture.}

TSVP

\page{53}
\begin{tikzcd}
  \widetilde{X}_2 \arrow[d] & \cdot \arrow[l] \\
  X_4 & X_8 \arrow[l]
\end{tikzcd}
\note{c'est le seul contenu de la page ; il reprend le diagramme de la p.~52. La source de la flèche horizontale du haut n'est pas écrite. Le $X$ de $X_8$ est écrit sur une première lettre \struck{\ill{}}.}

\page{54}
\[
  X = \frac{y^2 - x^2}{z^2}, \qquad Y = -\frac{2xy}{z^2}
\]
\[
  X + iY = \frac{1}{z^2}\bigl(y^2 - x^2 - 2ixy\bigr) = \frac{1}{z^2}\,(y - ix)^2
\]
\[
  (X + iY)^{n} = \frac{1}{z^{2n}}\,(y - ix)^{2n}
  = \frac{1}{z^{2n}}\bigl(P_n(x, y) + i\,Q_n(x, y)\bigr)
\]
\[
  \overset{?}{=}\; \frac{y'^2 - x'^2}{z'^2} \mp 2i\,\frac{x'y'}{z'^2}
\]
\note{il écrit $(y + ix)^2$ à la fin de la deuxième ligne, où le calcul donne $(y - ix)^2$ ; le signe est surchargé.}

$x', y', z'$ fonctions rat. des $x, y, z$ ?! P.ex. $z' = z^n$
\struck{$P_n(x, y) =$} \quad $\dfrac{B_n(x, y, z)}{Z^{n}}$ \quad et \quad $B_n$
\note{le $Z$ du dénominateur est lui-même \struck{biffé}.}
\begin{align*}
  B_n(x, y)^2 - A_n(x, y)^2 &= P_n(x, y) \\
  -2\, A_n(x, y)\, B_n(x, y) &= Q_n(x, y)
\end{align*}
\[
  X + iY = \Bigl(\frac{y - ix}{z}\Bigr)^2 = \Bigl(\frac{y}{z} - i\,\frac{x}{z}\Bigr)^2
\]
\[
  \Bigl(\frac{y}{z} - i\,\frac{x}{z}\Bigr)
  \qquad
  \Bigl(\frac{y}{z} - i\,\frac{x}{z}\Bigr)^{2n+1}
  \qquad
  (-i)^3 = -i\,(i^2) = i
\]
\[
  \Bigl(\frac{y}{z} - i\,\frac{x}{z}\Bigr)^3
  = \frac{1}{z^3}\,\bigl[\,y^3 - 3iy^2x - 3yx^2 + ix^3\,\bigr]
\]
\[
  = y^3 - 3yx^2 + i\,(x^3 - 3y^2x)
\]
\[
  = y\,(y^2 - 3x^2) + i\,x\,(x^2 - 3y^2)
\]
\note{les deux dernières lignes, sous une accolade, omettent le facteur $1/z^3$. Elles redonnent les $y'$ et $x'$ de la p.~48.}

\section{Automorphismes de groupoïdes}
\note{la p.~56, feuille rose, ne porte que les mots « autom. des groupoïdes », de sa main, en haut à droite : c'est une chemise pour les pp.~57--60.}

\page{57}
\emph{Automorphismes de groupoïdes.}

1) Soient $C$ \struck{\ill{}} un groupoïde, $C'$ une catégorie, \ill{} \ill{} :
\[
  \mathrm{Hom}(C, C') = \mathrm{Ob}\, \underline{\mathrm{Hom}}(C, C').
\]
On a
\[
  \mathrm{Hom}(C, C') \longrightarrow \mathrm{Hom}\bigl(\underbrace{\mathrm{Ob}\,C}_{C_0},\, \underbrace{\mathrm{Ob}\,C'}_{C'_0}\bigr)
\]
Soit
\[
  f_0 : C_0 \to C'_0 ,
\]
on \uncertain{cherche} les $f : C \to C'$ qui lui donnent naissance. On trouve :

\emph{Proposition} a) Pour qu'il existe un $f$ au-dessus de $f_0$, il faut et il suffit que pour $x, y \in C_0$ \uncertain{isomorphes} : c'est-à-dire \uncertain{dans la même} comp. connexe de $C$ (i.e. tels que $\mathrm{Hom}_C(x, y) \neq \emptyset$) on ait $f_0(x) \simeq f_0(y)$.
\struck{b) \ill{}, cette condition \ill{} \ill{}}
b) Pour toute composante connexe $C(i)$ de $C$, soit $a_i \in C(i)_0$, et pour tout $x \in C(i)_0$, soit
$u_x : a_i \xrightarrow{\ \sim\ } x$. Alors
\[
  \mathrm{Hom}_{f_0}(C, C') \xrightarrow{\ \sim\ }
  \prod_{x \in C_0} \mathrm{Isom}\bigl(f_0(a_{i(x)}), f_0(x)\bigr)
\]
\[
  \times \prod_{i} \mathrm{Hom}\bigl(\mathrm{Aut}(a_i), \mathrm{Aut}(f_0(a_i))\bigr)
\]
\note{sous le premier produit, un indice surchargé : $\struck{x \in C_0}$ puis \ill{}.}

\page{58}
\[
  f \longmapsto \Bigl(\bigl(\underbrace{f(u_x)}_{\ell_x}\bigr)_{x \in C_0},\
  \bigl(\underbrace{f_{a_i, a_i} : \mathrm{Aut}(a_i) \to \mathrm{Aut}(f_0(a_i))}_{\varphi_i}\bigr)\Bigr).
\]

c) Pour que $f$ soit fidèle, il faut et il suffit que $\forall i \in I$, $\varphi_i$ soit injectif.

Pour que $f$ soit pleinement fidèle, il faut et il suffit que
1°) pour \uncertain{$i \neq j$}, on ait $\mathrm{Hom}\bigl(f_0(a_i), f_0(a_j)\bigr) = \emptyset$ ;
2°) pour $\forall x, y \in C(i)_0$,
\[
  \mathrm{Isom}\bigl(f_0(x), f_0(y)\bigr) \xrightarrow{\ \sim\ } \mathrm{Hom}\bigl(f_0(x), f_0(y)\bigr) ;
\]
\note{le premier « Isom » est \uncertain{douteux}, surchargé sur une première écriture.}
3°) $\forall i \in I$, $\varphi_i$ est un iso.

\struck{\uncertain{détermine} $C$}
Pour que $f$ soit une équivalence, il faut et il suffit qu'il soit pl.\ fidèle et qu'on ait de plus
4°) Tout $x' \in C'_0$ est isomorphe à un des $f_0(a_i)$.

Pour que $f$ soit un iso, il faut et il suffit qu'on ait 1°, 2°, 3°) et de plus
4'°) $f_0$ bijective.

\page{59}
\note{le haut de la page, jusqu'à « $C_0 \xrightarrow{u_0} C'_0$ », est barré de quatre longs traits obliques ; il se lit néanmoins.}
\struck{$C = (C_0, C_1)$ groupoïde connexe}

\struck{$G = \mathrm{Aut}(C)$ groupe des automorphismes}

\struck{$G^0$ s-groupe des \ill{} de $G$ induisant l'identité sur $C_0$}
\struck{$1 \to G^0 \to G \to \mathfrak{S}_{C_0} \to 1$ (sous $G$ : \uncertain{surj})}

\struck{[Pour qu'une autoéquivalence \add{$u$} soit un iso, il faut et il suffit qu'elle soit bijective sur l'ensemble des objets de $C$, et induise un iso.}
\struck{\struck{sur les} $\mathrm{Aut}(a) \to \mathrm{Aut}(u(a))$ pour \struck{un} $x \in C_0$ ; d'autre part, \uncertain{toute} application $C_0 \xrightarrow{u_0} C'_0$}

\[
  1 \to \mathrm{Aut}^0(C) \to \mathrm{Aut}\,C \to \mathfrak{S}_{C} \to 1
\]
\struck{$\mathrm{Aut}^0(C) \simeq$}
\[
  1 \to \mathrm{Aut}^0(C, a) \to \mathrm{Aut}^0(C) \to \mathrm{Aut}(\pi_a) \to 1,
  \qquad \pi_a = \mathrm{Aut}_C(a)
\]
\struck{A} Par définition \struck{$\mathrm{Aut}^0(C, a) =$}
\[
  \mathrm{Aut}^0(C, a) = \bigl\{ f \in \mathrm{Aut}(C) \bigm| \mathrm{Ob} f = \mathrm{id},\
  \bigl(f_a : \mathrm{Aut}(a) \to \mathrm{Aut}(\underbrace{f(a)}_{=a})\bigr) = \mathrm{id} \bigr\}
\]
On considère, pour tout $x \in C_0$,
\[
  T_x = \mathrm{Isom}(a, x), \quad \pi_a\text{-torseur à droite}
\]
et son commutant \struck{\ill{}} $\mathrm{Aut}(x)$

\page{60}
\emph{Corollaire.} Pour que $f$ soit un iso, il faut et il suffit que
\begin{itemize}
  \item a) $f$ induise $\pi_0(f) : \pi_0(C) \xrightarrow{\ \sim\ } \pi_0(C')$
  \item b) $f$ induise $\mathrm{Ob} f = f_0 : \mathrm{Ob}(C) \xrightarrow{\ \sim\ } \mathrm{Ob}(C')$
  \item c) $\forall i \in I$, $f_{a_i, a_i} : \mathrm{Aut}(a_i) \xrightarrow{\ \sim\ } \mathrm{Aut}(f(a_i))$
\end{itemize}
\emph{NB} $f$ équivalence $\Longleftrightarrow$ a) et c) sont vrais.

\emph{Corollaire} Pour qu'une application $f_0 : C_0 \to C'_0$ provienne d'un iso $f : C \simeq C'$, il faut et il suffit que
\begin{itemize}
  \item a) $f_0$ \struck{est} bijective
  \item b) $f_0$ transforme la relation d'équivalence de $C_0$ (par comp.\ connexes) en celle de $C'_0$ (condition vide si $C, C'$ connexes).
  \item c) $\forall i \in I$, $\mathrm{Aut}\, f_0(a_i)$ est isomorphe à $\mathrm{Aut}(a_i)$.
\end{itemize}

\struck{\ill{}} Supposons $C = C'$. Alors on trouve
\[
  \mathrm{Aut}(C) \longrightarrow \mathfrak{S}_{C_0} \quad \text{\textit{surjective}}
\]
si $\mathrm{Aut}^0(C)$ est le noyau, on trouve dans
\note{la phrase se poursuit au-delà de ce lot.}

\end{document}
